\section{Conclusión}\label{sec:conclusion}
Las tres rondas muestran que el método más informativo depende del destino y del criterio aplicado. La comparación se apoya en medianas de IP únicas, cobertura y estabilidad, y conserva el costo de exploración.

Valparaíso: scamper: icmp, icmp-paris, udp, udp-paris (mediana 9 IP); ripe: ICMP, TCP, UDP (mediana 9 IP).

La Habana: scamper: icmp, icmp-paris (mediana 13 IP); ripe: TCP (mediana 14 IP).

Rusia: scamper: icmp, icmp-paris (mediana 12 IP); ripe: TCP (mediana 11 IP).

tracelb permite revelar interfaces adicionales en algunos modos y destinos, con distinto costo en sondas. Las diferencias entre Chile, Cuba y Rusia describen las rutas y horarios medidos; no deben interpretarse como una clasificación de velocidades o de censura.

Ida y vuelta no ofrecieron evidencia suficiente para afirmar simetría completa en las rutas incompletas. Las observaciones corresponden a interfaces; la identificación de routers físicos requiere corroboración adicional de alias. La respuesta selectiva sí se repitió en los casos documentados, pero no identifica por sí sola una regla de filtrado.

El aprendizaje principal es que la calidad de Path Discovery depende tanto del método como de la captura, la recepción local y la trazabilidad. Para repetir el estudio conviene preservar originales, contrastar parámetros efectivos, verificar la continuidad de la sesión y comparar pares próximos temporalmente. El diagnóstico TCP demuestra por qué una ausencia en la salida de una herramienta debe contrastarse con los paquetes antes de atribuirla a la red.

